% Part: many-valued-logic
% Chapter: syntax-and-semantics
% Section: semantic-notions

\documentclass[../../../include/open-logic-section]{subfiles}

\begin{document}

\olfileid[fr]{mvl}{syn}{sem}

\olsection{Notions sémantiques}

Supposons qu'une matrice détermine
une logique plurivalente~$\Log L$.
Nous pouvons alors définir pour
$\Log L$ les notions sémantiques usuelles.

\begin{defn}
\begin{enumerate}
\item Une formule~$!A$ est
  \emph{satisfaisable} s'il existe
  une valuation~$\pAssign{v}$ telle
  que $\pSat{v}{!A}$; elle est
  \emph{insatisfaisable} si aucune
  valuation ne la satisfait.
\item Une formule~$!A$ est une
  \emph{tautologie} si
  $\pSat{v}{!A}$ pour toute
  valuation~$v$.
\item Si~$\Gamma$ est un ensemble
  de formules, $\Gamma\Entails !A$
  (« $\Gamma$ entraîne~$!A$ »)
  si et seulement si toute valuation
  $\pAssign{v}$ qui satisfait
  $\Gamma$ satisfait aussi~$!A$.
\item Un ensemble de formules~$\Gamma$
  est \emph{satisfaisable} s'il existe
  une valuation~$\pAssign{v}$ telle
  que $\pSat{v}{\Gamma}$; sinon,
  il est \emph{insatisfaisable}.
\end{enumerate}
\end{defn}

Comme en logique classique, ces
notions vérifient certains faits :

\begin{prop}
\ollabel{prop:semanticalfacts}
\begin{enumerate}
\item $!A$ est une tautologie si et
  seulement si $\emptyset\Entails !A$;
\item si $\Gamma$ est satisfaisable,
  tout sous-ensemble fini de~$\Gamma$
  l'est aussi;
\item\ollabel{def:monotonicity}%
  \emph{Monotonie} : si
  $\Gamma\subseteq\Delta$ et
  $\Gamma\Entails !A$, alors
  $\Delta\Entails !A$;
\item\ollabel{def:Cut}%
  \emph{Transitivité} : si
  $\Gamma\Entails !A$ et
  $\Delta\cup\{!A\}\Entails !B$,
  alors $\Gamma\cup\Delta\Entails !B$.
\end{enumerate}
\end{prop}

\begin{proof}
Exercice.
\end{proof}

\begin{prob}
Démontrer la \olref[mvl][syn][sem]{prop:semanticalfacts}.
\end{prob}

En logique classique, la conséquence
est liée au conditionnel. Ainsi,
le \emph{modus ponens} est valide :
si $\Gamma\Entails !A$ et
$\Gamma\Entails !A\lif !B$,
alors $\Gamma\Entails !B$.
Le théorème sémantique de déduction
donne une autre relation :
$\Gamma\Entails !A\lif !B$
si et seulement si
$\Gamma\cup\{!A\}\Entails !B$.
Ces résultats ne valent
\emph{pas toujours} dans les
logiques plurivalentes : ils dépendent
de la fonction de vérité~$\tf{\lif}$.

\end{document}
